\documentclass[10pt]{article}
\usepackage[T1]{fontenc}
\usepackage[utf8]{inputenc}
\usepackage[french]{babel}
\usepackage{amsmath,amssymb,mathrsfs,mathpazo}
\usepackage{geometry}
\geometry{a4paper,margin=2cm}
\usepackage{graphicx,array,tabularx,tabularray,colortbl,multirow,multicol,booktabs,enumitem,pifont,ulem,fancybox,diagbox,cancel}
\usepackage{pstricks,pstricks-add,pst-plot,pst-tree,pst-node,pst-text,pst-func,pst-math,pst-bspline,pst-3dplot,pst-eucl,pst-solides3d}
\let\pscalculate\undefined
\usepackage{pst-bezier}
\usepackage{tikz,pgfplots,tkz-tab}
\usetikzlibrary{arrows,arrows.meta,patterns,calc,positioning,trees,decorations.markings,shapes.geometric,shapes.misc,shapes.arrows,intersections,decorations.pathmorphing,decorations.pathreplacing,backgrounds,shadows,babel}
\pgfplotsset{compat=1.18}
\usepgfplotslibrary{fillbetween}
\usepackage[np]{numprint}
\usepackage{listings,makecell,pas-tableur,twoopt}
\let\barre\undefined
\usepackage[pstricks]{tablvar}
\definecolor{darkgreen}{RGB}{0,100,0}
\providecommand{\red}{\color{red}}
\providecommand{\blue}{\color{blue}}
\providecommand{\green}{\color{green}}
\providecommand{\black}{\color{black}}
\providecommand{\white}{\color{white}}
\providecommand{\gray}{\color{gray}}
\providecommand{\cyan}{\color{cyan}}
\providecommand{\magenta}{\color{magenta}}
\providecommand{\yellow}{\color{yellow}}
\providecommand{\violet}{\color{violet}}
\providecommand{\orange}{\color{orange}}
\providecommand{\euro}{\texteuro}
\providecommand{\EUR}{\texteuro}
\providecommand{\ieme}{\textsuperscript{e}}
\providecommand{\degres}{\ensuremath{^\circ}}
\providecommand{\degree}{\ensuremath{^\circ}}
\providecommand{\vv}[1]{\overrightarrow{#1}}
\providecommand{\Coord}[1]{#1}
\providecommand{\decosix}{}
\providecommand{\starredbullet}{\ensuremath{\bullet}}
\pagestyle{empty}
\setlength{\parindent}{0pt}
\usepackage{adjustbox}\DeclareUnicodeCharacter{2236}{:}\DeclareUnicodeCharacter{2190}{\ensuremath{\leftarrow}}
\providecommand{\R}{\mathbb{R}}
\providecommand{\N}{\mathbb{N}}
\providecommand{\Z}{\mathbb{Z}}
\providecommand{\Q}{\mathbb{Q}}
\providecommand{\D}{\mathbb{D}}
\providecommand{\C}{\mathbb{C}}
\providecommand{\e}{{\mathrm{e}}}
\providecommand{\vect}[1]{\overrightarrow{#1}}
\providecommand{\vectt}[1]{\overrightarrow{\mathrm{#1}}}
\providecommand{\vvt}[1]{\overrightarrow{\mathrm{#1}}}
\providecommand{\np}[2][]{{#2~#1}}
\providecommand{\vv}[1]{{\overrightarrow{#1}}}
\renewcommand{\d}{{\mathrm{d}}}
\providecommand{\degres}{{^\circ}}
\providecommand{\degre}{{^\circ}}
\providecommand{\blue}{\color{blue}}
\providecommand{\red}{\color{red}}
\providecommand{\enskip}{\quad}
\providecommand{\llbracket}{\lbrack\!\lbrack}
\providecommand{\rrbracket}{\rbrack\!\rbrack}
\providecommand{\euro}{{\text{€}}}
\providecommand{\no}{{\text{n°}}}
\providecommand{\interval}[2]{{[#1;#2]}}
\providecommand{\intervFO}[2]{{[#1;#2[}}
\providecommand{\dotfill}{\ldots}

\providecommand{\ds}{\displaystyle}
\providecommand{\textding}[1]{#1}
\providecommand{\Oij}{\(\mathrm{O};\vec{\imath},\vec{\jmath}\)}
\providecommand{\Oijk}{\(\mathrm{O};\vec{\imath},\vec{\jmath},\vec{k}\)}
\providecommand{\Ouv}{\(\mathrm{O};\vec{u},\vec{v}\)}
\providecommand{\intervFF}[2]{\ensuremath{[#1;#2]}}
\renewcommand{\euro}{\ifmmode\text{\texteuro}\else\texteuro\fi}
\DeclareUnicodeCharacter{00B0}{\ifmmode^{\circ}\else\textdegree\fi}
\DeclareUnicodeCharacter{20AC}{\ifmmode\text{\texteuro}\else\texteuro\fi}
% Énoncés : Éducation nationale. Transcription et corrigés : APMEP, Denis Vergès et contributeurs. Découpage et classement : CoopMaths, Sébastien Lozano, Éric Elter et contributeurs. Adaptation de la présentation : Maths Explorer.
% CC BY-SA : https://coopmaths.fr/www/about/
% https://coopmaths.fr/alea/static/bac/2021/tex/bac_2021_06_metropole_10.tex
\begin{document}
\noindent\textit{APMEP / CoopMaths. CC BY-SA. Présentation adaptée par Maths Explorer.}\par\medskip
\smallskip

\begin{tabular}{|l|}\hline
Principaux domaines abordés :\\
Fonction exponentielle; dérivation; convexité\\ \hline
\end{tabular}

\begin{center}\textbf{Partie 1}\end{center}

On donne ci-dessous, dans le plan rapporté à un repère orthonormé, la courbe représentant la
fonction dérivée $f'$ d'une fonction $f$ dérivable sur $\R$.

À l'aide de cette courbe, conjecturer, en justifiant les réponses : 

\medskip

\begin{enumerate}
\item Le sens de variation de la fonction $f$ sur $\R$.
\item La convexité de la fonction $f$ sur $\R$.
\end{enumerate}

\begin{center}
\psset{unit=1cm,arrowsize=2pt 3}
\begin{pspicture*}(-3.25,-2)(5.25,4.5)
\psgrid[gridlabels=0pt,subgriddiv=1,gridwidth=0.15pt](-2,-1)(4,4.5)
\psaxes[linewidth=1.25pt](0,0)(-2,-1.25)(4,4.5)
\psaxes[linewidth=1.25pt]{->}(0,0)(1,1)
\psplot[plotpoints=2000,linewidth=1.25pt,linecolor=red]{-2}{4}{1 x add 2.71828 x exp div  neg}
\rput(1,-1.5){Courbe représentant la \textbf{dérivée} $f'$ de la fonction $f$.}
\end{pspicture*}
\end{center}

\smallskip

\begin{center}\textbf{Partie 2}\end{center}

\smallskip

On admet que la fonction $f$ mentionnée dans la Partie 1 est définie sur $\R$ par :

\[f(x) = (x + 2)\e^{-x}.\]

On note $\mathcal{C}$ la courbe représentative de $f$ dans un repère orthonormé \Oij.

On admet que la fonction $f$ est deux fois dérivable sur $\R$, et on note $f'$ et $f''$ les fonctions dérivées première et seconde de $f$ respectivement.

\medskip

\begin{enumerate}
\item Montrer que, pour tout nombre réel $x$,

\[f(x) = \dfrac{x}{\e^{x}}+ 2\e^{-x}.\]

En déduire la limite de $f$ en $+ \infty$.

Justifier que la courbe $\mathcal{C}$ admet une asymptote que l'on précisera.

On admet que $\displaystyle\lim_{x \to - \infty} f(x) = - \infty$.
\item 
	\begin{enumerate}
		\item Montrer que, pour tout nombre réel $x$,\, $f'(x) = (- x - 1)\e^{-x}$.
		\item Étudier les variations sur $\R$ de la fonction $f$ et dresser son tableau de variations.
		\item Montrer que l'équation $f(x) = 2$ admet une unique solution $\alpha$ sur l'intervalle $[-2~;~-1]$ dont on donnera une valeur approchée à $10^{-1}$ près.
	\end{enumerate}
\item Déterminer, pour tout nombre réel $x$, l'expression de $f''(x)$ et étudier la convexité de la fonction~$f$. 

Que représente pour la courbe $\mathcal{C}$ son point A d'abscisse $0$ ?
\end{enumerate}

\end{document}
